\clearpage
\section{Reproducibility details}
\label{app:repro}

\subsection{Configuration}

Both runs are launched from the same notebook with the same configuration block; the run identity
is encoded in the experiment name, which also determines the output directory, so two runs can
never share a folder. The tracked \texttt{run\_metadata.json} of each run records the full resolved
configuration. Diffing the two files yields exactly two substantive differences
(\texttt{lookahead\_k}: absent vs.\ 5; \texttt{lookahead\_sub\_batch\_size}: absent vs.\ 128) plus
the run's own name, adapter repository, two output paths, and start timestamp.

\begin{table}[h]
\centering
\small
\setlength{\tabcolsep}{3pt}
\resizebox{\columnwidth}{!}{%
\begin{tabular}{@{}ll@{}}
\toprule
Setting & Value \\
\midrule
Therapist model & \texttt{Llama-3.2-1B}, bf16 \\
LoRA & $r{=}16$, $\alpha{=}16$ \\
Learning rate & $1\mathrm{e}{-}5$ \\
Seed & 42 \\
Iterations & 10 \\
Epochs / iteration & 2 \\
Conversations / iteration & 96 \\
Group size $G$ & 8 \\
KL $\beta$ & 0.01 \\
GRPO temperature & 1.2 \\
GRPO loss variant & standard (\texttt{grpo}) \\
Policy updates per sampled batch & 1 \\
Train batch $\times$ accumulation & $64 \times 2$ \\
Trainer eval split & 0.05 of prompts \\
Min.\ context length (MCL) & 12 utterances \\
Session cap & 49 utterances after the opening \\
Max tokens / response & 200 \\
Therapist context & 2{,}048 tokens \\
Therapist / patient temperature & 0.9 / 0.7 \\
Patient + oracle & \texttt{gpt-4o-mini-2024-07-18} \\
Software & TRL 1.4.0, transformers 5.8.1, PEFT 0.19.1 \\
Hardware & one NVIDIA A100 (Google Colab) \\
Look-ahead $K$ & 0 ($\Kz$ run) / 5 ($\Kf$ run) \\
Optimizer steps, 10 iterations & $1{,}128$ ($\Kz$) / $1{,}070$ ($\Kf$) \\
\bottomrule
\end{tabular}}
\caption{Resolved configuration, identical across the two runs except the final two rows.
Source: each run's \texttt{run\_metadata.json}; the step counts are read from the per-step
training artifacts (\S\ref{app:repro-cost}).}
\label{tab:config}
\end{table}

\subsection{Instruments}
\label{app:instruments}

\begin{table*}[t]
\centering
\small
\setlength{\tabcolsep}{6pt}
\begin{tabular}{@{}llll@{}}
\toprule
Instrument & Source & Outputs per conversation & Reported here \\
\midrule
Q1 & \citet{yosef2024assessing} & 5 items, 1--5 & mean \\
Q2 & \citet{yosef2024assessing} & 17 items, 1--5 & mean \\
WAI-SR & \citet{hatcher2006waisr} & 12 items, 1--5 & mean \\
CSQ-8 & \citet{larsen1979csq} & 8 items, 1--4 & mean \\
MI-SAT & adapted, this work & 6 items, 1--5 & mean \\
MITI & \citet{moyers2016miti} & 4 globals, 1--5; 7 behaviour counts & mean of the globals \\
PCT & this work & 3 globals, 1--5; 3 utterance counts & CT / (CT + ST) \\
MICI & this work & 1 global, 1--5; 6 behaviour counts & acts per therapist turn \\
\midrule
utterance coder & this work & one code per therapist and per patient utterance & \S\ref{sec:therapist}--\ref{sec:patient} \\
\bottomrule
\end{tabular}
\caption{The eight instruments and the utterance coder. \QQ\ is the mean of the Q1 and Q2 means. CT
and ST are the patient's change-talk and sustain-talk utterance counts. MICI is lower-is-better;
its per-channel counts and rates are the material of \S\ref{sec:therapist}. The utterance coder is
described in \S\ref{app:coder}.}
\label{tab:instruments}
\end{table*}

Table~\ref{tab:instruments} lists the eight instruments. Each is administered as one API call per
conversation with a JSON-schema-constrained response, the fixed instructions and item texts first
and the transcript last. For Q1 and Q2 the judge is addressed as ``a professional motivational
interview therapist'' and told to rate each item from the patient's point of view on a 1--5 scale,
every item carrying a one-paragraph explanation with an example of a good and of a bad therapist
response (Q1: overall satisfaction with the chat and with its content, motivation facilitated,
whether anything was learned, and its relevance to everyday life; Q2: seventeen relational items
such as ``the therapist seemed to understand me'', ``the therapist did not judge me'', and ``the
therapist made me feel like he cared about me''). WAI-SR, CSQ-8 and MI-SAT are scored from the
patient's perspective with their item wording and response anchors. The three coders are addressed
as a certified MITI coder. The MITI coder returns four global ratings (cultivating change talk,
softening sustain talk, partnership, empathy) and assigns exactly one of seven behaviour codes to
each therapist utterance (giving information, persuade, question, simple reflection, complex
reflection, affirm, seeking collaboration). The PCT coder rates the patient's expressed importance,
confidence and readiness to change and classifies every patient utterance as change talk, sustain
talk or neutral. The MICI coder rates overall severity and counts six MI-inconsistent behaviours
across the therapist's turns, any number per utterance: confront, advise without permission, warn,
direct, judge/label, and over-praise, the last defined to the judge as ``effusive, generic, or
excessive praise NOT tied to a specific client strength/effort (`you're amazing', `I'm so proud of
you'). A genuine, specific MI affirmation does NOT count here.'' The per-turn rates in
\S\ref{sec:therapist} divide these counts by the number of therapist utterances, which the prompt
states.

\subsection{The utterance coder}
\label{app:coder}

The utterance coder is administered as one further call per conversation under each judge, with
the same rubric-first layout. The judge is addressed as an MI process coder trained on the MITI
4.2.1 and MISC manuals and told to assign exactly one code to every therapist utterance, its
dominant function (the function that occupies most of the utterance; a turn that praises at
length and ends in one question is coded as praise), and exactly one code to every patient
utterance, in transcript order. The therapist codes are: open question (invites elaboration);
closed question; simple reflection (repeats or rephrases, adding little or no meaning); complex
reflection (adds substantial meaning or emphasis: an inferred feeling, a metaphor, a double-sided
reflection); affirmation, in the MITI~4.2.1 sense of a comment on a specific client strength,
ability, intention or effort; non-specific praise (cheerleading or effusive reassurance not tied to
a specific strength or effort, including praise for a step the patient has not taken); giving
information; persuasion (arguing, lecturing, advising toward change, warning, moralising, with or
without permission); seeking collaboration or emphasising autonomy; confront / direct / judge; and
other (greeting, filler, logistics, degenerate text). The patient codes are change talk (any
statement favouring change), sustain talk (any statement favouring the status quo) and neutral,
with an utterance containing both coded by the direction that dominates. The transcript is
numbered by role (``[THERAPIST \#k]'', ``[PATIENT \#k]'') so that the arrays align, their
lengths are pinned to the utterance counts in the response schema, a response of the wrong length
is rejected and re-requested, and the scripted opener is pinned to open question and excluded
from every rate and transition. Every one of the $2 \times 11 \times 96 = 2{,}112$ evaluation
conversations is coded under both judges; three held-out rows that came back one code short on
the batch path were re-scored on the live path.

All process quantities are recomputed from the stored code sequences. Alternation is strict and
the therapist speaks first, so therapist turn $i$ is answered by patient turn $i$ and answers
patient turn $i-1$. The therapist's replies (Table~\ref{tab:process}) condition on the patient's
preceding code: of the therapist turns that follow a patient's sustain talk, for example, the share
that are persuasion. Change-talk persistence is, within a conversation, the share of the patient's
change-talk utterances whose next patient utterance is change talk again (and, after sustain talk,
the share followed by change talk); every such rate is computed per conversation and averaged over
the conversations in which its condition occurs. As a check of the
coder against the session-level instruments of the same judge, the per-utterance codes summed
per conversation are compared with the MITI behaviour counts and the PCT change-talk counts:
across the 21 model states, the patient-side counts agree closely (Spearman $\rho$ pooled within
state: change talk $0.88$ / $0.92$, sustain talk $0.90$ / $0.94$ under the training oracle / the
held-out judge), while the therapist-side counts agree only partly (training oracle
$0.02$--$0.43$ over the seven MITI codes; held-out judge $0.73$ for questions, $0.68$ for giving
information, $0.55$ for persuasion, $0.10$--$0.30$ for reflections, affirmations and seeking
collaboration). The two instruments count different things: the MITI coder counts every function a
long turn performs, the utterance coder its dominant one, which is why the MITI coder reports
$5.5$ questions per conversation where the utterance coder reports $1.7$ question-turns under the
training oracle, and why \S\ref{sec:therapist} counts open-question \emph{turns}.

\subsection{The simulated patient and the therapist prompt}
\label{app:prompts}

The patient's system prompt is assembled from the six persona attributes, quoted here verbatim
including spelling: ``You are speaking with a motivational interviewing counselor therapist, and
you are the patient in this conversation. Your name is [James / Emma], and you are [27 / 61] years
old [male / female]. [cooperation clause]. [problem clause]. [history clause].'' followed by an
instruction to avoid repetition and the rule for ending the session (write ``SESSION ENDED'' if
the therapist is wrapping up, or if the patient is satisfied and believes they gained enough
knowledge). The three cooperation clauses are: high (``Your level of cooperation is very high'');
low (``Your level of cooperation is very low, you should be very suspicious about the advices you
get from the therapist. In general you don't believe in therapy and you should respond
adverserially to the therapist. You should cooperate only if the therpist said something that
makes sense to you. You went to therapy only because someone forced you to do so''); and
low-then-high (``In the beginning of the session, you are less cooperative, you should be very
suspicious about the advices you get from the therapist. In general you don't believe in therapy
and you should respond adverserially to the therapist. You went to therapy only because someone
forced you to do so. But as the session progresses, you become more cooperative and more
motivated to change''). The problem clause describes smoking (``You have been
smoking for [a few months / many years], and it has become a daily habit. You are increasingly
concerned about the impact of smoking on your health'') or obesity (``You have been struggling
with obesity for [a few months / many years]. Your weight is negatively impacting your health. You
have high blood pressure and experience joint pain''), and the history clause states either that
the patient has never tried to change or that they have tried many times and relapsed. The
therapist's system prompt, applied through the chat template in both runs and inside the
look-ahead rollout, is ``You are a motivational interviewing counselor named David. You partner
with the patient to understand his problems. You are empathetic towards him and help the patient
explore their ambivalence regarding behavioral change. You are non-judgmental while encouraging
the patient to change,'' followed by the same anti-repetition instruction and its own
session-ending rule. Every conversation opens with the same scripted therapist line, omitted from
Appendix~\ref{app:example}.

\subsection{The lexical over-praise marker}
\label{app:marker}

The keyword marker of \S\ref{sec:therapist} is one case-insensitive regular expression matched against
each therapist utterance; a turn counts if it contains any of ten patterns: ``I'm so proud'',
``proud of you'', ``inspiration to me'', ``you got this'', ``beautiful'', ``beacon'',
``shining'', ``warrior'', ``hero of your'', and ``you are a light'' or ``you are a beacon''. It
exists as a direction check on the oracle's over-praise count and is deliberately crude
(``beautiful'' matches any use of the word); it is reported as the share of therapist turns
matched, never as a count of acts, and its absolute level is not a measurement. Its agreement with
the coded praise share (Figure~\ref{fig:process}a) is what \S\ref{sec:therapist} relies on.

\subsection{Anti-degeneracy}

The base 1B model will occasionally continue past its own turn and write the patient's reply. All
decode sites stop at the chat-template markers and truncate any completion at the first such
marker; a completion that is empty after cleaning ends the conversation, and a degenerate
completion is floored to the minimum reward. This applies identically in both runs, including
inside the $K{=}5$ rollout.

\subsection{Evaluation and statistics}

Model states are scored offline into a score archive keyed by judge, replicate, instrument and
model, one record per conversation; an interrupted scoring run resumes by skipping completed work.
The reported draw is replicate 0 for each judge. The Base was drawn twice, once at the start of
each run, with identical settings; the paper pools the two draws into one Base of 192
conversations, and a persona-paired contrast against the Base uses the mean of that persona's two
conversations.

All contrasts are persona-paired. The 96 personas are reshuffled every iteration, so pairing
recovers the persona identity by replaying the shuffle; pairing on file order would join unrelated
conversations across iterations. Tests
are Wilcoxon signed-rank; effect sizes are Cohen's \dz; intervals are 2{,}000-resample percentile
bootstraps at a fixed seed, so a re-render reproduces the same bounds. Holm correction is applied
within each family (instruments within one matched contrast, or iterations within one sweep) and
is not pooled across families.

\subsection{Cost accounting behind the disclosures}
\label{app:repro-cost}

The API-call counts quoted in the Limitations section are exact counts read from the per-iteration
generation logs (one record per oracle call and per patient call, summed over the ten training
iterations of each run), and the optimizer-step counts are the number of per-step training
artifacts each run wrote. The GPU-hour totals in the Limitations and the Ethics Statement and the
per-step wall-clock multiplier are reconstructed from the modification times of those artifacts, because
the run metadata's elapsed-time fields are per-process and undercount any iteration that crashed
and resumed, in the worst case here by nearly $2\times$. An inter-step gap outside
$(0, 3600\,\mathrm{s})$ is treated as a resume gap and imputed at the phase median; imputation
counts are recorded per run. The first two iterations of the $\Kf$ run ran at a smaller look-ahead
sub-batch, so the per-step multiplier is quoted over the settled iterations 3--10.

\subsection{Artifacts}

Every number in this paper is traceable to a tracked table in the analysis-results tree of the
code release, through a claims ledger released with it. Every figure is drawn by scripts in the
paper folder from that tree's tracked tables (Figure~\ref{fig:schematic}, the schematic, from no
data), at the width it is included at. The
transcripts in Table~\ref{tab:excerpt} and Appendix~\ref{app:example} are reproduced verbatim
from the stored conversation files, and the selection scripts are released with the paper.
